%%%%

%%%   Самарский государственный технический университет

%%%   Кафедра прикладной математики и информатики

%%%

%%%   Преамбула для статьи в журнал "Вестник Самарского государственного

%%%   технического университета. Серия: «Физико-математические науки»"

%%%   Ориентирована под MikTEx 2.4 for Win32

%%%

%%%   Хотя сам журнал собирается под GNU/Linux на дистрибутиве TeXLive



\documentclass[11pt,a4paper,twoside]{article}



\usepackage[cp1251]{inputenc}

\usepackage[T2A]{fontenc}

\usepackage[english,russian]{babel}

\usepackage{samgtubib}

\usepackage[dvips]{graphicx}

\usepackage[multiple]{footmisc}



\usepackage{samgtu}

\renewcommand{\VestnikYear}{2009} % Год издания Вестника

\renewcommand{\VestnikNumber}{2\,(19)} % Номер Вестника

\renewcommand{\VestnikMonth}{Ноябрь} % Месяц выхода Вестника

\renewcommand{\VestnikData}{01.11.09} % Дата подписания Вестника в печать

\renewcommand{\VestnikUPL}{26,64} % Усл. печ. л.

\renewcommand{\VestnikUIL}{25,73} % Уч.-изд. л.

\renewcommand{\VestnikRegNumber}{479/09} % Регистрационный номер

\renewcommand{\VestnikOrder}{994} % Номер заказа











%\begin{document}





\renewcommand{\firstpage}{108} %Начальная страница статьи для выноса в колонтитул


\renewcommand{\lastpage}{112} %Конечная страница статьи для выноса в колонтитул





%\Partition{Математическое моделирование}





\udk{519.246}





\title{Определение динамических характеристик нелинейных диссипативных 


систем на~основе линейно параметрической дискретной модели 


амплитудно"=частотной характеристики} % Полное название статьи


      {Определение динамических характеристик нелинейных диссипативных систем \dots} % Название статьи для колонтитула





\engtitle{Nonlinear dissipative systems dynamical characteristics determination based on linear--parametrical discrete model of~amplitude--frequency characteristic}





\author{Д.\,Н.~Попова}{П\,о\,п\,о\,в\,а~Д.\,Н.}





\engauthor{D.\,N.~Popova}





\date{01.11.2007} % Дата поступления статьи в редакцию.





\thanks{}





\address{\samgtu}





\engaddress{Samara State Technical University, Samara, Russia}% Название организации на англ.. языке





\email{dashek@mail.ru}





\workabstract{Рассматривается линейно параметрическая дискретная модель амплитудно"=частотной характеристики механической системы с~нелинейными диссипативными силами в~форме стохастического разностного уравнения. Предлагается эффективный метод определения динамических характеристик, в~основе которого лежит численная итерационная процедура среднеквадратичного оценивания коэффициентов линейно параметрической дискретной модели.}





\engabstract{Considering the mechanical system linear--parametrical discrete model of AFC with dissipative nonlinear forces in form of accidental difference equation, this work offers effective method of dynamical characteristics determination based on numeral iterative procedure of LPD model mean-square coefficient evaluation}





\maketitle








Качественные изменения в машиностроении неразрывно связаны с~оценкой 


технического состояния механических систем. Одними из основных методов при 


диагностике технического состояния механических систем являются методы, 


использующие анализ изменения динамических характеристик системы в процессе 


её эксплуатации, прочностных или других испытаний. Проблема повышения 


надёжности и~точности прогнозирования долговечности и~ресурса предъявляет 


высокие требования к~точности оценок динамических характеристик механических 


систем.





Известные способы оценки динамических характеристик диссипативных 


механических систем и~математические модели, лежащие в~их основе, не 


позволяют в~полной мере использовать преимущества развития современных 


технологий при обработке экспериментальных данных. Эта проблема может быть 


устранена при помощи методов, в~основе которых лежат параметрические модели 


временных рядов, широкое используемые в~спектральном анализе. Однако анализ 


известных параметрических моделей, в частности, моделей 


авторегрессии "--- скользящего среднего, показал их неэффективность при оценке 


динамических характеристик нелинейных диссипативных систем, так как в~них не 


описаны связи между динамическими характеристиками системы и~параметрами 


модели.





Наиболее эффективным и перспективным методом определения динамических 


характеристик нелинейной диссипативной механической системы является метод, 


в~основе которого лежит линейно параметрическая дискретная модель (ЛПДМ) 


свободных колебаний системы в форме стохастического разностного уравнения~[1]. 


Согласно этому методу с помощью итерационной процедуры вычисляются 


среднеквадратичные оценки коэффициентов ЛПДМ, через которые по известным 


формулам находятся динамические характеристики системы, в~том числе, степень 


нелинейности диссипативной силы. Такой способ принципиально отличается от 


традиционных методов нахождения динамических характеристик диссипативной 


системы~[2], он позволяет обеспечить высокую точность оценивания за счёт 


эффективного использования современных средств вычислений и~их математического обеспечения.





В данной работе рассматривается метод определения динамических характеристик 


нелинейной диссипативной системы, в~основе которого лежит итерационная 


процедура среднеквадратичного оценивания коэффициентов ЛПДМ в~форме 


стохастического разностного уравнения, построенного для амплитудно"=частотной 


характеристики.





Дифференциальное уравнение, описывающее вынужденные колебания системы с~


диссипативными силами, пропорциональной $n$-ной степени скорости движения, 


может быть представлено в~виде~[3]:


\[


m{y}''(t) + b{y}'(t)\left| {{y}'(t)} \right|^{n - 1} + cy(t) = P(t),


\]


где $m$ "--- масса системы; $c$ "--- коэффициенты жёсткости; $b$ "--- коэффициент в~нелинейной диссипативной 


силе; $n$ "--- показатель нелинейности диссипативной силы; $P(t)$ "--- возмущающее воздействие на систему.





При гармоническом воздействии $P(t) = P_\ast \sin \omega t$ дифференциальное 


уравнение, описывающее вынужденные гармонические колебания, может быть 


переписано в~виде


\[


{y}''(t) + \frac{b}{m}b{y}'(t)\left| {{y}'(t)} \right|^{n - 1} + p^2y(t) = 


P_0 \sin \omega t,


\]


где $p^2 = \frac{c}{m}$ "--- резонансная частота; $\omega $ и $P_0 = \frac{P_\ast }{m}$ 


"--- частота и~приведённая к~массе системы амплитуда внешней гармонической силы. 


Решением этого уравнения является установившиеся колебания \mbox{$y(t) = a\cos (\omega t + \varphi )$}, зависимость амплитуды которых (в~относительных к~$y_{\text{ст}} $ единицах, где \mbox{$y_{\text{ст}} = \frac{P_0 }{p^2})$} от частоты возбуждения 


(амплитудно"=частотная характеристика  (АЧХ) системы) описывается функцией вида~[3]:


\begin{equation}


\label{popova:eq1}


\tilde {a}(\omega ) = \frac{1}{\sqrt {\left( {1 - \frac{\omega ^2}{p^2}} 


\right)^2 + \left( {\frac{bS_n }{\pi c}a^{n - 1}\omega ^n} \right)^2} }.


\end{equation}








Для систем с диссипативными силами, пропорциональными $n$-ой степени скорости 


движения (частотно"=зависимым трением), зависимость декремента колебаний от 


амплитуды имеет вид


\begin{equation}


\label{popova:eq2}


\delta (a) = \frac{bp^n}{c}S_n a^{n - 1}.


\end{equation}





Тогда амплитудно"=частотную характеристику (\ref{popova:eq1}) систем с~диссипативными 


силами, пропорциональными $n$-ной степени скорости движения, с~учётом (\ref{popova:eq2}) можно 


представить в~виде


\begin{equation}


\label{popova:eq3}


\tilde {a}(\omega ) = \frac{1}{\sqrt {\left( {1 - \frac{\omega ^2}{p^2}} 


\right)^2 + \left( {\frac{\delta }{\pi }} \right)^2\left( {\frac{\omega 


^2}{p^2}} \right)^n} }.


\end{equation}


Представляя выражение (\ref{popova:eq3}) в~виде $$\tilde {a}(\omega ) = \frac{1}{\sqrt 


{\left( {1 - \frac{\omega ^2}{p^2}} \right)^2 + \frac{\delta ^2}{\pi 


^2}\left( {1 + \left[ {\frac{\omega ^2}{p^2} - 1} \right]} \right)^n} }$$ и~используя разложение 


в~степенной ряд: $$\left( {1 + \left[ 


{\frac{\omega ^2}{p^2} - 1} \right]} \right)^n \approx 1 + n\left( 


{\frac{\omega ^2}{p^2} - 1} \right),$$ получаем $$\tilde {a}(\omega ) = 


\frac{1}{\sqrt {\left( {1 - \frac{\omega ^2}{p^2}} \right)^2 + \frac{\delta 


^2}{\pi ^2}\left( {1 + n\left( {\frac{\omega ^2}{p^2} - 1} \right)} \right)} 


}.$$





Полагая $\omega = \Delta \omega k$ ($k = 0,\,1,\,2,\, \dots$), где $\Delta \omega 


$ "--- период равномерной дискретизации АЧХ, получаем её дискретной аналог: 


$$\tilde {a}_k = \frac{1}{\sqrt {\left( {1 - \frac{\Delta \omega ^2k^2}{p^2}} \right)^2 + \frac{\delta ^2}{\pi ^2}\left( {1 + n\left( {\frac{\Delta \omega ^2k^2}{p^2} - 1} \right)} \right)} }.$$ 


Отсюда после алгебраических преобразований получаем $$\tilde {a}_k ^2\left[ {\frac{\Delta \omega 


^4}{p^4}k^4 + \left( {\frac{\delta ^2}{\pi ^2}n - 2} \right)\frac{\Delta 


\omega ^2}{p^2}k^2 + \frac{\delta ^2}{\pi ^2}\left( {1 - n} \right) + 1} 


\right] = 1.$$ Применяя к~обеим частям данного выражения Z"=преобразование, 


получаем $$Z\left\{ {\tilde {a}_k ^2\left[ {\frac{\Delta \omega ^4}{p^4}k^4 + 


\left( {\frac{\delta ^2}{\pi ^2}n - 2} \right)\frac{\Delta \omega 


^2}{p^2}k^2 + \frac{\delta ^2}{\pi ^2}\left( {1 - n} \right) + 1} \right]} 


\right\} = \frac{1}{1 - z^{ - 1}}.$$ Введём обозначение $\mu _0 = 


\frac{\Delta \omega ^4}{p^4}$, $\mu _1 = \left( {\frac{\delta ^2}{\pi ^2}n - 


2} \right)\frac{\Delta \omega ^2}{p^2}$, $\mu _2 = \frac{\delta ^2}{\pi 


^2}\left( {1 - n} \right) + 1$, отсюда получаем, что $$\left( {1 - z^{ - 1}} 


\right)Z\left\{ {\tilde {a}_k^2 \left( {\mu _0 k^4 + \mu _1 k^2 + \mu _2 } 


\right)} \right\} = 1.$$ Возвращаясь в~пространство оригиналов, используя 


первую теорему смещения: $$z^{ - r}Z\left\{ {\tilde {a}_k } \right\} = \tilde 


{a}_{k - r},\quad  r = 0,\,1,$$ а также полагая, что $\tilde {a}_{k - r} = 0$ 


при $k - r < 0$, получаем разностное уравнение вида $$\left( {\mu _0 k^4 + \mu 


_1 k^2 + \mu _2 } \right)\tilde {a}_k^2 - \left( {\mu _0 \left( {k - 1} 


\right)^4 + \mu _1 \left( {k - 1} \right)^2 + \mu _2 } \right)\tilde {a}_{k 


- 1}^2 = \delta _k,\quad  k = 0,\,1,\,2,\,\dots ,$$ где 


$\delta _k = \left\{ {\begin{array}{lll}


 1,&\text{при} & k = 0, \\ 


 0,&\text{при} & k \ne 0.


 \end{array}} \right.$ "--- символ Кронекера.





При $k = 0$ получаем, что


\begin{equation}


\label{popova:eq4}


\tilde {a}_0^2 = \frac{\pi ^2}{\delta ^2(1 - n) + \pi ^2}.


\end{equation}


При $k \geqslant 1$ получаем разностное уравнение вида


\begin{equation}


\label{popova:eq5}


\lambda _0 \left( {\tilde {a}_k^2 k^4 - \tilde {a}_{k - 1}^2 \left( {k - 1} 


\right)^4} \right) + \lambda _1 \left( {\tilde {a}_k^2 k^2 - \tilde {a}_{k - 


1}^2 \left( {k - 1} \right)^2} \right) = \tilde {a}_{k - 1}^2 - \tilde 


{a}_k^2 ,


\end{equation}


где $\lambda _0 = \frac{\mu _0 }{\mu _2 } = \frac{\Delta \omega ^4\pi 


^2}{p^2\left( {\delta ^2 - \delta ^2n + \pi ^2} \right)}$, $\lambda _1 = 


\frac{\mu _1 }{\mu _2 } = \frac{\left( {\delta ^2n - 2\pi ^2} \right)\Delta 


\omega ^2}{p^2\left( {\delta ^2 - \delta ^2n + \pi ^2} \right)}$.





Результаты наблюдений $a_k $ содержат аддитивную случайную помеху 


$\varepsilon _k $: $a_k = \tilde {a}_k + \varepsilon _k $, где $\tilde {a}_k 


$ "--- мгновенные значения АЧХ, соответствующие формуле~(\ref{popova:eq3}).





С учётом случайной помехи в~результатах измерений, линейно параметрическую 


дискретную модель АЧХ можно представить в~форме стохастических разностных 


уравнений:


\begin{equation}


\label{popova:eq6}


\left\{ \begin{array}{l}


 a_0^2 = \lambda _2 - \varepsilon _0^2 + 2\varepsilon _0 a_0 , \\ 


 a_{k - 1}^2 - a_k^2 = \lambda _0 \left( {a_k^2 k^4 - a_{k - 1}^2 (k - 1)^4} 


\right) + \lambda _1 \left( {a_k^2 k^2 - a_{k - 1}^2 (k - 1)^2} \right) + 


\eta _k , \\ 


 \eta _k = \left[ {\lambda _0 k^4 + \lambda _1 k^2 + 1} \right]\left( 


{\varepsilon _k^2 - 2\varepsilon _k a_k } \right) - \left[ {\lambda _0 (k - 1)^4 + \lambda _1 (k - 1)^2 + 1} \right]\times\\ 


\qquad \qquad   \qquad \qquad  \qquad \qquad \times\left( 


{\varepsilon _{k - 1}^2 - 2\varepsilon _{k - 1} a_{k - 1} } \right),


\quad k = 1,\,2,\,\dots,\, N - 1.


 \end{array} \right.


\end{equation}








Пренебрегая величинами второго порядка малости относительно $\varepsilon _k 


$, систему (\ref{popova:eq6}) можно представить в~виде:


\begin{equation}


\label{popova:eq7}


\left\{ {\begin{array}{l}


 a_0^2 = \lambda _2 + \eta _0 , \\ 


 a_{k - 1}^2 - a_k^2 = \lambda _0 \left( {a_k^2 k^4 - a_{k - 1}^2 (k - 1)^4} 


\right) + \lambda _1 \left( {a_k^2 k^2 - a_{k - 1}^2 (k - 1)^2} \right) + 


\eta _k , \\ 


 \eta _0 = 2\varepsilon _0 a_0 , \\ 


 \eta _k = - 2\varepsilon _k a_k \left[ {\lambda _0 k^4 + \lambda _1 k^2 + 


1} \right] + 2\varepsilon _{k - 1} a_{k - 1} \left[ {\lambda _0 (k - 1)^4 + 


\lambda _1 (k - 1)^2 + 1} \right],


 \end{array}} \right.


\end{equation}


где $k = 1,\,2,\,\dots,\, N - 1.$





Коэффициенты в модели (\ref{popova:eq7}) связаны с~параметрами АЧХ следующими соотношениями:


\begin{equation}


\label{popova:eq8}


\lambda _0 = \frac{\Delta \omega ^4\pi ^2}{p^2\left( {\delta ^2 - \delta ^2n 


+ \pi ^2} \right)},


\quad


\lambda _1 = \frac{\left( {\delta ^2n - 2\pi ^2} \right)\Delta \omega 


^2}{p^2\left( {\delta ^2 - \delta ^2n + \pi ^2} \right)},


\quad


\lambda _2 = \frac{\pi ^2}{\delta ^2 - \delta ^2n + \pi ^2}.


\end{equation}





В матричной форме ЛПДМ (\ref{popova:eq7}), описывающая последовательность результатов измерений мгновенных 


значений АЧХ, имеет вид


\begin{equation}


\label{popova:eq9}


\left\{ {\begin{array}{l}


 b = F\lambda + \eta , \\ 


 \eta = P_{a,\lambda } \varepsilon . 


 \end{array}} \right.


\end{equation}


Здесь $\lambda = (\lambda _0 ,\,\lambda _1 ,\,\lambda _2 ,\,\lambda _3 )^{\rm T}$ 


"--- вектор неизвестных коэффициентов ЛПДМ; 


$\varepsilon = (\varepsilon _0,\,\varepsilon _1,\, \dots,\, \varepsilon _{N - 1} )^{\rm T}$ "--- $N$-мерный вектор 


случайной помехи в~результатах наблюдений; 


$\eta = (\eta _0 ,\,\eta _1 ,\, \dots,\,\eta _{N - 1} )^{\rm T}$ "--- $N$-мерный вектор эквивалентного возмущения в~стохастическом разностном уравнении; $b = (a_0^2,\, a_0^2 - a_1^2,\, \dots,$ $ \, a_{N - 2}^2 - a_{N - 1}^2 )^{\rm T}$ "--- $N$-мерный вектор правой части; $F$ "--- матрица регрессоров размера $N\times 4$, столбцы которой описываются 


следующими формулами:


\begin{gather*}


f_{i1} = (0,\,a_1^2 ,\,16a_2^2 - a_1^2 ,\, \dots,\,(N - 1)^4a_{N - 1}^2 - (N - 2)^4a_{N - 2}^2 )^{\rm T},


\\


f_{i2} = (0,\,a_1^2 ,\,4a_2^2 - a_1^2 ,\, \dots,\,(N - 1)^2a_{N - 1}^2 - (N - 2)^2a_{N - 2}^2 )^{\rm T},


\\


f_{i3} = (1,\,0,\,\dots,\,0)^{\rm T},\quad i = 1,\, 2,\, \dots,\,  N.


\end{gather*}











Матрица  $P$ размера $N\times N$ в~стохастическом разностном уравнении 


эквивалентного возмущения "--- нижняя треугольная, ленточная, двухдиагональная. 


Первый столбец матрицы $P$ имеет вид 


$p_{i1} = (2a_0 ,\,2a_0 ,\,0,\, \dots,\,0)^{\rm T}$ ($i = 1,\, 2,\, \dots,\,  N$). Остальные 


столбцы матрицы $P$ ($j = 2,\, 3,\, \dots,\,  N$) описываются 


формулами


\begin{equation}


\label{popova:eq10}


p_{ij} = \left\{ \begin{array}{cll}


 0, & \text{если} & i < j\;\text{и}\;i > j + 1, \\ 


- 2a_{i - 1} \left( {\lambda _0 (i - 1)^4 + \lambda _1 (i - 1)^2 + 1} 


\right), & \text{если} & i = j, \\ 


- p_{i - 1,j} , & \text{если} & i = j + 1.


 \end{array} \right.


\end{equation}











Построенная в форме стохастических разностных уравнений, линейно 


параметрическая дискретная модель служит основой нового, помехозащищённого 


метода определения динамических характеристик систем с диссипативными силами 


по их АЧХ. Этот метод включает следующие основные этапы:





1) формирование выборки


$a_k$ ($k = 0,\, 1,\, \dots,\, N - 1$), результатов измерений с~шагом $\Delta \omega$


амплитудно-частотной характеристики системы, где $N$ "---  объём выборки;





2) вычисление по приведённым выше формулам элементов матрицы $F$ и~вектора $b$;





3) среднеквадратичное оценивание коэффициентов $\lambda_i$ ($i=0,\,1$) из первого стохастического разностного уравнения системы~(10); при этом наиболее эффективным является итерационный метод среднеквадратичного оценивания коэффициентов линейно параметрической дискретной модели~[1]; в~соответствии с~этим методом вначале вычисляются МНК"=оценки 


$\hat\lambda^{(0)}= (F^{\rm T}F)^{-1}F^{\rm T}b$, на основе которых формируется первое приближение матрицы $P_\lambda$; затем первое уравнение в~системе (10)~преобразуется к~виду $P_{\hat\lambda}^{-1}b = P_{\hat\lambda}^{-1} F\lambda +\varepsilon$ и вновь вычисляются МНК"=оценки, но уже для преобразованного уравнения, при этом процедура уточнения повторяется несколько раз;





4) вычисление на основе среднеквадратичных оценок $\hat\lambda_i$ параметров системы с~диссипативными силами по формулам


\begin{equation}


\label{popova:eq11}


\hat {p} = \sqrt[4]{\frac{\hat {\lambda }_2 }{\hat {\lambda }_0 }}\Delta 


\omega ,


\quad


\hat {\delta } = \pi \sqrt {1 + \frac{\sqrt {\hat {\lambda }_0 } + \hat 


{\lambda }_1 \sqrt {\hat {\lambda }_2 } }{\sqrt {\hat {\lambda }_0 \hat 


{\lambda }_2 } }} ,


\quad


\hat {n} = \frac{\lambda _1 \sqrt {\lambda _0 \lambda _2 } + 2\lambda _0 


\lambda _2 }{\lambda _1 \sqrt {\lambda _0 \lambda _2 } + \lambda _0 \lambda 


_2 + \lambda _0 }.


\end{equation}








\begin{wrapfigure}{O}{.53\textwidth}


\centering


\includegraphics{./00Pic/popova.eps}


\parbox[b]{.5\textwidth}{\small Зависимости относительной погрешности вычисления декремента колебаний от величины случайной помехи в результатах наблюдений при использовании методов определение по ширине пика (кривая 1) и на основе линейно параметрической дискретной модели (кривая~2)}


\vspace{-5mm}


\end{wrapfigure}








Проведены численно"=аналитические исследования по сравнению данного метода 


с~существующим методом определения декремента по кривой резонанса. Численный 


эксперимент проводился следующим образом. Генерировалась выборка $a_k$  


($k = 0,\, 1,\, \dots,\,N - 1$) дискретных значений функции, описывающей амплитудно"=частотную 


характеристику в~системе с~диссипативными силами, пропорциональными $n$-степени скорости движения, 


с~параметрами $\delta _0 = 0{,}05$, $p = 2\pi $, $n = 2$. Период дискретизации 


выбирался таким образом, что при $T = \frac{2\pi }{\omega _p } = 1$ 


отношение $\frac{\Delta \omega }{T} = 0{,}1$. Объём выборки составил $N =100$. 


В~отсчёты $\tilde {a}_k $, соответствующие точным значениям функции, добавлялась 


аддитивная помеха $\varepsilon _k$, мощность которой изменялась  в~диапазоне от 0~до~3\,{\%}. 


Для каждой точки эксперимента результаты вычислений относительной погрешности усреднялись по $M = 100$ реализациям.





На рисунке представлены зависимости относительной погрешности 


вычислений декремента колебаний от мощности случайной помехи. Точки~1 соответствуют методу определения декремента по кривой резонанса, точки~2 "--- рассмотренному методу.





Из результатов исследования видно, что погрешность оценивания декремента 


колебаний по предложенному методу на порядок выше точности оценивания 


декремента колебаний по кривой резонанса. Таким образом, применение метода 


определения динамических характеристик на основе ЛПДМ в форме 


стохастического разностного уравнения, построенного для АЧХ диссипативной механической системы, увеличивает точность 


вычисления декремента по сравнению с~методом определения декремента по кривой резонанса.





\begin{thebibliography}{3}


\item \textit{Зотеев, В.\,Е.} Определение динамических характеристик нелинейных диссипативных систем на 


основе стохастического разностного уравнения [Текст]~/ В.\,Е.~Зотеев, Д.\,Н.~Попова~/\!/ Вестн. Сам.


гос. техн. ун-та. Сер.: Физ.-мат. науки. "---  2006. "---  \No~42. "---  C.~162--168. "--- ISBN 5--7964--0815--1.





\item \textit{Писаренко, Г.\,С.} Методы определения характеристик колебаний упругих систем [Текст]~/ Г.\,С.~Писаренко, В.\,А.~Матвеев, А.\,П.~Яковлев. "--- Киев: Наук. думка, 1976. "--- 88~с.





\item \textit{Пановко, Г.\,Я.} Основы прикладной теории колебаний и удара [Текст]~/ Г.\,Я.~Пановко. "--- Л.: Машиностроение, 1976. "--- 320~с.


  


\end{thebibliography}





\vspace{8mm}


\lfoot[]{}


  \rfoot[]{}





\makeengtitle





%\newpage


%\makecontents





%\end{document}


